\corrigesChapitre{4}{Équations polynomiales}
\begin{corrige}{1}{Un second degré}
\begin{questions}
\item $z=\pm\mathrm i$.
\item $z=\pm2\mathrm i$.
\item $(z-2\mathrm i)(z+2\mathrm i)$.
\end{questions}
\end{corrige}
\begin{corrige}{2}{Discriminant}
\begin{questions}
\item $1\pm2\mathrm i$.
\item Somme $2$, produit $5$.
\item $(z-z_1)(z-z_2)=z^2-(z_1+z_2)z+z_1z_2=z^2-2z+5$.
\end{questions}
\end{corrige}
\begin{corrige}{3}{Différence de puissances}
\begin{questions}
\item $(z-1)(z^2+z+1)$.
\item $(z-2)(z^3+2z^2+4z+8)$.
\item $3$, un degré de moins que le polynôme initial.
\end{questions}
\end{corrige}
\begin{corrige}{4}{Vérifier une racine}
\begin{questions}
\item $1-3+4-2=0$.
\item Le développement donne $z^3-3z^2+4z-2$.
\item $z=1$ ou $z^2-2z+2=0$ ; les solutions sont $1,1+\mathrm i,1-\mathrm i$.
\end{questions}
\end{corrige}
\begin{corrige}{5}{Troisième degré}
\begin{questions}
\item $8-8+2-2=0$.
\item $P=(z-2)(z^2+1)$ par développement ou identification.
\item $2,\mathrm i,-\mathrm i$.
\end{questions}
\end{corrige}
\begin{corrige}{6}{Coefficient dominant}
\begin{questions}
\item $\Delta=16-80=-64$.
\item $z=(-4\pm8\mathrm i)/4=-1\pm2\mathrm i$.
\item $2(z+1-2\mathrm i)(z+1+2\mathrm i)$ ; le facteur $2$ ne doit pas disparaître.
\end{questions}
\end{corrige}
\begin{corrige}{7}{Racines conjuguées}
\begin{questions}
\item $2-\mathrm i$, par conjugaison.
\item $(z-2-\mathrm i)(z-2+\mathrm i)=z^2-4z+5$.
\item $(z+1)(z^2-4z+5)=z^3-3z^2+z+5$.
\end{questions}
\end{corrige}
\begin{corrige}{8}{Une équation bicarrée}
\begin{questions}
\item $w^2+5w+4=(w+1)(w+4)=0$, donc $w=-1$ ou $w=-4$.
\item $z=\pm\mathrm i$ ou $z=\pm2\mathrm i$.
\item Il y a quatre racines distinctes, ce qui atteint la borne du degré $4$.
\end{questions}
\end{corrige}
\begin{corrige}{9}{Racine double}
\begin{questions}
\item $P(z)=z^3-3z+2$.
\item Deux : $1$ et $-2$ ; $1$ est répété dans la factorisation.
\item Le théorème donne au plus trois racines distinctes, et non nécessairement trois.
\end{questions}
\end{corrige}
\begin{corrige}{10}{Une identité de polynômes}
\begin{questions}
\item La différence est un polynôme de degré au plus $2$ possédant trois racines distinctes.
\item Si elle était non nulle, elle aurait au plus deux racines. Elle est donc le polynôme nul.
\item Vérifier l’égalité en trois points distincts suffit lorsqu’on connaît effectivement la borne de degré.
\end{questions}
\end{corrige}
\begin{corrige}{11}{Carrée d’un complexe, guidée}
\begin{questions}
\item $x^2-y^2=3$ et $2xy=4$.
\item $x^2+y^2=5$ ; addition et soustraction donnent $x^2=4,y^2=1$.
\item La contrainte $xy=2$ impose les mêmes signes : $w=2+\mathrm i$ ou $w=-2-\mathrm i$.
\end{questions}
\end{corrige}
\begin{corrige}{12}{Preuve par le degré}
\begin{questions}
\item Sa valeur ne change jamais et reste différente de zéro.
\item Pour une racine $b\ne a$, $0=(b-a)Q(b)$ avec $b-a\ne0$, donc $Q(b)=0$.
\item Si $P$ a degré $n$, $Q$ a degré $n-1$ ; il possède au plus $n-1$ racines distinctes. En ajoutant $a$, on obtient la borne $n$.
\end{questions}
\end{corrige}
